\documentclass[11pt]{article}
\usepackage[margin=1in]{geometry}
\usepackage{amsmath,amssymb,hyperref}
\title{A counterexample to Conjecture 00000008656}
\author{gaochengzhecpu}
\date{}
\begin{document}
\maketitle

The assertion that a minimized word is unique in its automorphism orbit is
false.  It fails in the free group of rank $2$, so no rank-one convention is
involved.  Refuting this asserted conjunct suffices to refute the stated
conjecture; no assertion about the other complexity clauses is needed.

\paragraph{Counterexample.}
Let $F=F(a,b)$ be the free group on two generators, with length measured in
the basis $\{a,b\}$.  Consider the reduced words
\[
 u=a,\qquad v=a^{-1}.
\]
The substitutions $a\mapsto a^{-1}$ and $b\mapsto b^{-1}$ extend to an
automorphism $\sigma$ of $F$.  Applying these substitutions twice gives the
identity.  This map substitutes inverses of the basis letters in their
original order; it is a group automorphism, not the map reversing products
that takes an arbitrary group element to its inverse.
In particular, $\sigma(u)=v$, so $u,v$ belong to the same automorphism orbit.

Both have length $1$.  Every automorphism is injective and fixes the identity.
Since $a\ne1$, every element in its automorphism orbit is nonidentity.  A word
of length $0$ represents only the identity.  Therefore every word representing
any element of the orbit has length at least $1$.  This proves minimality
against all automorphisms and all words, with no search cutoff.

Finally the homomorphism $F\to\mathbb Z$ sending both generators to $1$
sends $u$ to $1$ and $v$ to $-1$.  Thus $u\ne v$ even as group elements.
The orbit consequently contains at least two distinct minimized words.

\paragraph{Scope of the conclusion.}
This is the literal uniqueness assertion in both supplied versions of the
source: a unique minimum representative in an orbit.  The source supplies
no additional identification under signed basis permutations and no
lexicographic tie-breaking rule.  A different assertion about a chosen
canonical representative is not refuted here.  The two words are also
cyclically reduced.  For background on the standard problem of multiple
minimum-length words in one automorphism orbit, see Myasnikov and Shpilrain,
\emph{Automorphic orbits in free groups}, Section 1,
\href{https://shpilrain.ccny.cuny.edu/orbit.pdf}{author's full text}.
The elementary argument above is self-contained.

\newpage
\paragraph{Formal construction and complete quantifiers.}
The Lean file imports only \texttt{Std}.  For an arbitrary basis type $A$, it
constructs words as all finite lists of signed basis letters and constructs
$F(A)$ as their quotient by the congruence generated by
$xx^{-1}=1$.  It proves that concatenation and inverse words descend to this
quotient and verifies every group axiom.  The theorem
\texttt{free\_universal\_property} then proves existence and uniqueness of a
homomorphism to every group for every assignment of the basis.  This
establishes that the object used is the free group, not a finite surrogate or
an abelian replacement.

The type \texttt{Automorphism} consists of multiplication- and identity-
preserving maps with two-sided inverse functions.  The file constructs the
signed basis permutation on the quotient and proves that it is an
involutive automorphism.  The orbit predicate quantifies over this full type
of automorphisms.  The predicate \texttt{MinimalWord} compares a word with
every finite word representing every element of that orbit.  This stronger
comparison includes all reduced representatives and therefore establishes
the usual reduced-word minimality.  Both displayed words are also formally
checked to be reduced.

The integer-valued separating homomorphism is obtained from the proved
universal construction.  The main theorem
\begin{center}
\texttt{Conjecture8656.conjecture8656\_counterexample}
\end{center}
proves that the two length-one words are minimal, that their quotient group
elements are distinct, and that uniqueness fails.  Its only reported axioms
are the standard \texttt{propext} and \texttt{Quot.sound}.  There are no
\texttt{sorry}, \texttt{admit}, \texttt{native\_decide}, or custom axioms.
\end{document}
